\documentclass[10pt,a4paper]{amsart}
\usepackage[margin=19mm]{geometry}
\usepackage{amsmath,amssymb,mathtools}
\usepackage[scr=rsfs,cal=euler]{mathalfa}
\usepackage{xfrac}
\usepackage[unicode,hidelinks,bookmarks=false]{hyperref}
\newcommand{\ie}{i.e.\ }
\newcommand{\cf}{cf.\ }
\newcommand{\resp}{respectively\ }
\newcommand{\Subsection}{\subsection}
\providecommand{\nobreakdash}{\nobreak\hbox{-}}
\setlength{\emergencystretch}{2em}
\multlinegap0pt
\usepackage{bm}
\usepackage{bbm}
\usepackage{dsfont}
\usepackage{xspace}
\usepackage{cancel}
\usepackage{mathtools}
\usepackage{amsthm}
\makeatletter
\@ifundefined{lemma}{\newtheorem{lemma}{Lemma}}{}
\makeatother
\title[Counting Unlabeled Chordal Graphs by Equivariant Evaporation]{Counting Unlabeled Chordal Graphs by Equivariant Evaporation}
\author{Matthew Sun}
\date{Source: arXiv:2607.02613v1 (CC BY 4.0)}
\begin{document}
\maketitle

\noindent\emph{Source and licence.} Counting Unlabeled Chordal Graphs by Equivariant Evaporation. Version arXiv:2607.02613v1 (CC BY 4.0), distributed under the Creative Commons Attribution 4.0 International licence, as recorded in the arXiv licence metadata for this exact version and shown on the abstract page https://arxiv.org/abs/2607.02613v1. Full rights notice: licenses/arxiv-2607.02613-CC-BY-4.0.txt.\par\smallskip

\noindent\textit{Editorial scope.} Counted unit: the Lemma whose source label is \texttt{lem:percost} in the article (source0.tex lines 453--459, source label \texttt{lem:percost}), delivered together with the article's own complete original proof of it (source0.tex lines 461--486, one \texttt{proof} environment of 26 lines). No mathematics is written, repaired or re-derived: statement and proof are the article's own text line for line. Required context retained uncounted: eq:gf (lines 315--319). Every definition, display and label of those lines is retained unchanged. Omitted from this package and not redistributed: 1--314, 320--452, 460--460, 487--736. Nothing inside a retained excerpt is omitted or altered: every retained body is delivered exactly as the article's lines are written, so no omission receipt of any kind accompanies this package. Citations: the retained excerpts carry no citation command at all, so no bibliography entry of the article is transcribed and no reference is redistributed. Reference adaptation: none was needed and none was made; every reference inside the delivered unit resolves inside it, so no unresolved reference or citation is produced. Printed slips kept literal: no printed word, symbol, hypothesis or number of the counted statement or of its proof was altered. Layout and class: the article's own class is \texttt{article} with packages including geometry, amsmath, amssymb, amsthm, mathtools, booktabs, enumitem, hyperref; the wrapper is the lane's pinned \texttt{amsart} editorial wrapper at 10pt on A4, which restates the theorem-like environments the retained text needs and adds the article's own macro definitions that the retained text needs (none), with extra wrapper packages (bm, bbm, dsfont, xspace, cancel, mathtools). The delivered numbering is the unit's own. Assets: no retained span contains a figure environment, a graphics-inclusion command, a PDF-page inclusion, a table or a TikZ picture; no image file is copied into this package, the package declares no asset dependency and no image file is redistributed with it. Licence: version arXiv:2607.02613v1 of this article is distributed under the Creative Commons Attribution 4.0 International licence, as recorded in the arXiv licence metadata for this exact version and shown on the abstract page https://arxiv.org/abs/2607.02613v1. A full rights notice is delivered at licenses/arxiv-2607.02613-CC-BY-4.0.txt. Characters: every retained excerpt is pure ASCII, so the delivered engine, XeLaTeX with its default Latin Modern text font, reports no missing character.
\begin{equation}\label{eq:gf}
\Gamma_d(x)\;=\;\prod_{i=1}^{r}\big(1+v_{\tau_i}\big)^{\gcd(d,s_i)\,x_i},
\qquad
\Gamma_d^{\mathrm{cov}}(x)\;=\;\prod_{i=1}^{r}\Big(\big(1+v_{\tau_i}\big)^{\gcd(d,s_i)}-1\Big)^{x_i}.
\end{equation}

\begin{lemma}\label{lem:percost}
Let $r=r(\lambda)$ be the number of distinct part sizes of $\lambda$. The dynamic program for
$\lambda$, including its divisor-bundle sub-instances, runs in time
\[
T(\lambda)\ \le\ n^{O(1)}\,P(\lambda)^{O(1)}\,(n+1)^{\,O(r)}.
\]
\end{lemma}

\begin{proof}
\emph{States.} Each function is memoized over an index $t\le n$ and between one and four tuples
($x,k,z$, and $l$ for the $f$-functions), each bounded coordinatewise by the multiplicities, so
taking at most $P(\lambda)$ values; there are $O\!\big(n\,P(\lambda)^4\big)$ states per function.

\emph{Block generating functions.} A recurrence's bundle term builds a multivariate generating
function~\eqref{eq:gf} in $r'\le r$ block variables of total degree at most the number of blocks
$\sum_i\gcd(d,s_i)x_i\le n$. Such a polynomial has at most $\prod_\tau(\beta_d(x)_\tau+1)\le(n+1)^{r}$
monomials, and is assembled in $(n+1)^{O(r)}$ ring operations. (This corrects a naive
$n^{O(1)}$ estimate: the monomial count is exponential in the number of distinct block sub-types.)

\emph{Work per state.} Besides the generating function, a recurrence sums over a sub-multiset of $k$
($\le P(\lambda)$ terms; $\le P(\lambda)^2$ for the absorption recurrence) and over divisors
(at most $d(c)=n^{o(1)}\le n$), and evaluates the sub-world $g_1$ at each generating-function monomial (a memo
lookup). All counts are non-negative integers bounded by the number of labeled chordal graphs on $n$
vertices, $\le 2^{\binom n2}$, hence of $O(n^2)$ bits, so each arithmetic operation costs
$\mathrm{poly}(n)$. Work per state is therefore $n^{O(1)}\,P(\lambda)^{2}\,(n+1)^{r}$.

\emph{Sub-instances.} A bundle term calls a sub-instance whose multiplicity vector has
$P(\lambda')\le P(\lambda)$ (the map $s\mapsto s/\gcd(d,s)$ only merges parts, and
$(a{+}1)(b{+}1)\ge a{+}b{+}1$) and at most $r$ distinct sizes; the recursion depth is $\le\log_2 n$
(each step at least halves the largest size). The number of distinct reachable sub-instances is at
most $\prod_i d(s_i)\le 2^{O(r)}$; charging each by the same bound multiplies the total by $2^{O(r)}$,
which is absorbed into the $(n+1)^{O(r)}$ factor (as $(n+1)^{O(r)}\ge 2^{O(r)}$), leaving the stated
form.
\end{proof}

\end{document}
